% !TEX program = pdflatex
% THEME 1 — Difficile — Corrigé
\documentclass[11pt,a4paper]{article}
\usepackage[utf8]{inputenc}
\usepackage[T1]{fontenc}
\usepackage{lmodern}
\usepackage[french]{babel}
\usepackage{amsmath,amssymb,mathtools}
\usepackage{icomma}
\usepackage{siunitx}
\sisetup{output-decimal-marker={,},per-mode=symbol}
\DeclareSIUnit{\atm}{atm}
\DeclareSIUnit{\bar}{bar}
\DeclareSIUnit{\mmHg}{mmHg}
\usepackage[version=4]{mhchem}
\usepackage{xcolor}
\usepackage{colortbl}
\usepackage[most]{tcolorbox}
\usepackage{enumitem}
\usepackage{array}
\usepackage{booktabs}
\usepackage{fancyhdr}
\usepackage[a4paper,margin=2.1cm,headheight=26pt,headsep=10pt]{geometry}

\definecolor{primary}{RGB}{150,98,162}
\definecolor{primarydark}{RGB}{92,52,110}
\definecolor{excol}{RGB}{0,135,90}
\definecolor{attncol}{RGB}{200,40,55}

\newtcolorbox{rep}[1][]{enhanced,breakable,colback=excol!5,colframe=excol,
  boxrule=0.8pt,arc=2pt,left=3mm,right=3mm,top=2mm,bottom=2mm,
  fonttitle=\bfseries,coltitle=white,title={#1}}

\pagestyle{fancy}\fancyhf{}
\fancyhead[L]{\small\textcolor{primarydark}{\textbf{Fiche 7} — Loi des gaz}}
\fancyhead[R]{\small\textcolor{primarydark}{\textsc{Thème 1 · Difficile · Corrigé}}}
\fancyfoot[C]{\small\thepage}
\renewcommand{\headrulewidth}{0.8pt}
\renewcommand{\headrule}{\hbox to\headwidth{\color{primary}\leaders\hrule height \headrulewidth\hfill}}
\setlength{\parindent}{0pt}\setlength{\parskip}{3pt}

\newcounter{exo}
\newenvironment{cor}[1][]{\refstepcounter{exo}\medskip
  \noindent{\color{primarydark}\bfseries Corrigé \theexo.}\ \textit{#1}\par\nopagebreak\smallskip}{\par}
\newcommand{\rf}[1]{\textbf{\textcolor{attncol}{#1}}}

\begin{document}

\begin{center}
{\Large\bfseries\color{primarydark} Thème 1 — Niveau Difficile — Corrigé}
\end{center}
\textcolor{primary}{\hrule height 1pt}
\medskip

\begin{cor}[Cylindre à piston, étape par étape]
\textbf{(a) Étape A (isotherme).} $T$ constante $\Rightarrow$ \textbf{Boyle} : $p_1V_1=p_AV_A$.
\[ V_A=\frac{p_1V_1}{p_A}=\frac{\qty{1.0}{\atm}\times\qty{24}{\litre}}{\qty{2.0}{\atm}}=\rf{\qty{12}{\litre}}. \]
La pression double, le volume est divisé par $2$.

\textbf{(b) Étape B (isobare).} $p$ constante $\Rightarrow$ \textbf{Charles} : $\dfrac{V_A}{T_1}=\dfrac{V_B}{T_2}$,
avec $T_1=\qty{300}{\kelvin}$ et $T_2=\qty{400}{\kelvin}$.
\[ V_B=V_A\frac{T_2}{T_1}=\qty{12}{\litre}\times\frac{400}{300}=\rf{\qty{16}{\litre}}. \]

\textbf{(c) Loi combinée directe} de l'état 1 à l'état 3 :
\[ \frac{p_1V_1}{T_1}=\frac{p_3V_3}{T_3}\;\Rightarrow\;
   V_3=V_1\times\frac{p_1}{p_3}\times\frac{T_3}{T_1}
     =24\times\frac{1{,}0}{2{,}0}\times\frac{400}{300}=24\times0{,}5\times1{,}333=\rf{\qty{16}{\litre}}. \]
On retrouve exactement $V_B$ : \rf{cohérent}. Le chemin suivi (une ou deux étapes) n'a pas d'importance,
seuls comptent les états initial et final.

\textbf{(d)} On calcule $\dfrac{pV}{T}$ dans chaque état (en $\si{\atm\litre\per\kelvin}$) :
\[ \text{état 1 : }\frac{1\times24}{300}=0{,}08 \;;\quad
   \text{état 2 : }\frac{2\times12}{300}=0{,}08 \;;\quad
   \text{état 3 : }\frac{2\times16}{400}=0{,}08. \]
La quantité $\dfrac{pV}{T}$ est \rf{la même dans les trois états}. C'est normal : d'après $pV=nRT$, elle
vaut $nR$, qui ne dépend que de la quantité de gaz (constante ici). Ce sera le point de départ du
Thème 2.
\end{cor}

\begin{cor}[Mesure en unités SI]
\textbf{(a)} La pression, le volume et la température \emph{changent tous les trois} ; seule la quantité
$n$ reste constante. On utilise donc la \textbf{loi combinée} $\dfrac{p_1V_1}{T_1}=\dfrac{p_2V_2}{T_2}$.

\textbf{(b)} On isole $p_2$ :
\[ p_2=p_1\times\frac{V_1}{V_2}\times\frac{T_2}{T_1}. \]

\textbf{(c)} Application numérique :
\[ p_2=\qty{9.6e5}{\pascal}\times\frac{0{,}035}{0{,}0035}\times\frac{440}{220}
     =\qty{9.6e5}{\pascal}\times10\times2=\rf{\qty{1.92e7}{\pascal}}. \]
Aucune conversion n'est nécessaire car \emph{tout est déjà en unités SI} ($p$ en \si{\pascal}, $V$ en
\si{\meter\cubed}, $T$ en \si{\kelvin}) : dans la loi combinée, les unités de $V$ et de $T$ se simplifient
entre les deux membres, et $p_2$ ressort dans la même unité que $p_1$.

\textbf{(d)} La pression est passée de $9{,}6\times10^5$ à $1{,}92\times10^7$ : elle a \rf{augmenté d'un
facteur $20$} ($\times10$ par la compression, $\times2$ par le chauffage).
\end{cor}

\begin{cor}[Doubler la pression, et le zéro absolu]
Volume constant $\Rightarrow$ loi \textbf{isochore} $\dfrac{p_1}{T_1}=\dfrac{p_2}{T_2}$. On convertit :
$T_1=0+273=\qty{273}{\kelvin}$.

\textbf{(a)} Doubler la pression signifie $p_2=2p_1$, donc
\[ T_2=T_1\frac{p_2}{p_1}=273\times2=\qty{546}{\kelvin}
   \;\Rightarrow\; \theta_2=546-273=\rf{\qty{273}{\degreeCelsius}}. \]

\textbf{(b)} L'étudiant a doublé la température \emph{en degrés Celsius}. Or c'est la température
\rf{absolue (en kelvins)} qui est proportionnelle à la pression. Doubler $p$ exige de doubler $T$ en
kelvins ($273\to\qty{546}{\kelvin}$), ce qui correspond à $\qty{273}{\degreeCelsius}$ et non à
$\qty{0}{\degreeCelsius}$. À $\qty{0}{\degreeCelsius}$ le gaz est en réalité déjà à $\qty{273}{\kelvin}$,
pas à « zéro température ».

\textbf{(c)} La pression s'annule quand $T=\qty{0}{\kelvin}$ (car $p\propto T$). En degrés Celsius :
\[ \theta=0-273=\rf{\qty{-273}{\degreeCelsius}}. \]
C'est le \rf{zéro absolu} : la température la plus basse possible, où l'agitation des molécules — et donc
la pression qu'elles exercent — cesse.
\end{cor}

\end{document}
